\begin{frame}[fragile]{FAQ: 에폭이 많은 이유, $W_{hh}$ 미분}
  \begin{block}{Q. 왜 이렇게 많은 에폭이 필요했나요?}
    일부러 아주 단순한 구조(은닉 8)와 \textbf{기본 경사하강법}
    (Adam이 아니라)으로 \textbf{손으로 따라가기 쉽게} 만들었기
    때문 --- 실전에서는 Adam(Week 9의 SGD를 개선한 옵티마이저)
    + 더 큰 은닉 크기로 훨씬 적은 반복에 수렴.
    \kq{중요한 것은 몇 에폭이 걸렸는가가 아니라, 손실이 실제로
    0에 가깝게(12.09 $\to$ 0.00086) 줄었고 생성 문자열이 정답과
    일치했다는 사실 자체.}
  \end{block}
  \vskip0.6em
  \begin{block}{Q. $W_{hh}$가 공유되니까 BPTT에서 ``한 번만'' 미분하면
  안 되나요?}
    아니다 --- \textbf{미분은 각 시점에서 각각} 해야 하고, 그
    결과를 같은 $W_{hh}$의 그래디언트 텐서에 \textbf{더해야} 함
    (\verb|+=|).
    \begin{itemize}
      \item 공유된 파라미터의 그래디언트 = 계산 그래프에서
        \textbf{그 파라미터가 몇 번 쓰였는지만큼}의 기여 합
      \item CNN의 필터(Week 10)도 같은 원리: 하나의 필터가 이미지
        위에 $n^2$번 쓰이면 그 필터의 그래디언트는 $n^2$개 위치의
        \textbf{기여 합}
    \end{itemize}
  \end{block}
\end{frame}
